% ============================================================
% 第 5 章 案例分析 + 第 6 章 讨论 + 第 7 章 有效性威胁 + 第 8 章 结论
% ============================================================
\section{案例分析}\label{sec:case}

为理解量化结果背后的文本机制，本节选取一组典型的同源对比进行细读：V4.1-Flash 生成、T1（函数式编程语言）任务、第 1 次重复。该对由 V4-Pro 评审，\baseline{} 正文以 9/7（整体/原创）胜出。两篇输出的对比如图~\ref{fig:case}所示（节选）。

\begin{figure}[htbp]
\centering
\begin{minipage}[t]{0.485\textwidth}
\centering
\fcolorbox{gray!60}{gray!6}{\begin{minipage}{0.94\textwidth}
\small
\textbf{\baseline{}：Aster（星芒）}\par\vspace{0.3em}
``下面给出一份\textbf{可落地的}函数式编程语言设计草案。语言名为 Aster（星芒）：默认纯函数、强类型、显式效应、结构化并发、布局语法。\par
……\par
``Aster 是一门：默认纯函数式的语言；严格求值、默认不可变；使用代数数据类型 + 模式匹配；使用类型与效应行显式描述副作用；内建结构化并发；\textbf{小核心、大标准库，适合脚本、服务端和 WASM}。''\par
……\par
随后依次给出语法、类型系统（效应行）、并发模型、\textbf{实现路线图}与迁移策略等章节。
\end{minipage}}
\end{minipage}\hfill
\begin{minipage}[t]{0.485\textwidth}
\centering
\fcolorbox{gray!60}{orange!8}{\begin{minipage}{0.94\textwidth}
\small
\textbf{\wordpairs{}（清洗后）：潮织（Tideweave）}\par\vspace{0.3em}
``\textbf{潮织}是一种纯函数式、静态类型、以\textbf{效应光谱、潮汐并发、织纹模块、线性契约}为核心的实验性函数式编程语言。\par
……\par
``副作用被建模为效应光谱，用处理器进行折射与消解；结构化并发模型 tide 像潮汐一样汇聚、取消、回流结果；模块、类型和模式匹配采用织纹组合方式……''\par
……\par
正文围绕上述意象组织，术语新颖，但工程细节（实现路线、迁移策略）明显更简略。
\end{minipage}}
\end{minipage}
\caption{T1 任务第 1 次重复的两篇输出节选（V4.1-Flash 生成，V4-Pro 评审）。评审给分：Aster 整体 9 / 原创 7；潮织整体 7 / 原创 9。}
\label{fig:case}
\end{figure}

该案例浓缩了全局模式。\baseline{} 产物（Aster）从一开始就围绕“可落地”组织：非目标清单、类型与效应系统、实现路线图等章节一应俱全，评委理由为“结构更清晰、细节更具体，提供了明确的语法、类型系统、效应处理、并发模型和实现路线图，可落地的完整度更高”。\wordpairs{} 产物（潮织）则表现出鲜明的\textbf{意象驱动}特征：词对中的“潮·汐”“线·结”“织·梭”“光·谱”被系统性地转译为语言机制命名（tide 并发、织纹模块、效应光谱），评委承认其“隐喻和术语原创性更强”（原创性 9 对 7），但同时指出“在若干设计细节和可操作性上不如（\baseline{}）清晰”（可操作性 6 对 9）。这组九分制细节说明词对提示的作用通道确实存在——它把模型的搜索空间拉向了训练分布中更少出现的意象组合——但该收益以牺牲工程严谨性为代价，在多维度加权后净效应为负。

对 240 条评委理由的定性归纳显示，\baseline{} 的获胜理由高频出现“结构（更完整/清晰）”“细节（更具体/充分）”“可落地/可操作”“MVP（边界明确）”“路线图”等词；\wordpairs{} 的 15 次胜出理由则几乎全部集中于“（概念）更完整自洽”“原创整合度”“结构更丰富的分类学”等表述，且多数出现在 T5（人格模型）任务上。例如一条 T5 的评委理由写道：“A（\wordpairs{}）的模型结构更完整：八枢以四轴组织，九相由三动机核$\times$三调节流生成，并补充顺流、逆流、跨核代偿、影子枢…… 逻辑闭环与可操作性更强”。这与定量发现（T5 是唯一高于随机水平的任务）互为印证。

\section{讨论}\label{sec:discussion}

\subsection{词对提示为何整体失效：三条机制}

\textbf{机制一：指令目标稀释与预算挤占。}词对前缀实际上向模型同时下达了两个目标（先产出 200 组词、再完成任务），而两者共享同一生成预算与同一注意力资源。实验观测到三重挤占证据：(1) \wordpairs{} 条件完成 token 为 \baseline{} 的 $1.60\sim2.02$ 倍，说明词对环节本身消耗了大量生成预算；(2) 清洗后正文长度反而只有 \baseline{} 正文的 44\%--71\%，说明正文的绝对展开量下降；(3) 受损最重的维度恰是对“展开程度”最敏感的可操作性与具体性。对经过指令微调的模型而言，“先做 A 再做 B”的复合指令天然诱导对 A 的过度投入\cite{ouyang2022instructgpt}——词对是逐条枚举型内容，枚举的边际成本低而模型有完成清单的强倾向，正文于是成为预算的剩余部分。

\textbf{机制二：联想语境的自干扰。}词对列表一旦生成，就成为正文的直接上文。Shi 等人证明模型极易被无关上下文分散\cite{shi2023distracted}，而本研究中的“无关”并非人工注入而是模型自生成——这带来一个微妙的自增强回路：模型倾向于让自己的正文“呼应”刚写下的词对（潮织案例中的意象转译即为证据），哪怕这种呼应并不服务于任务质量。当联想词与任务高度相关时该回路可能有益；但 200 组无约束联想的期望相关性很低，多数词对成为正文生成时需要绕开的噪声。同时，正文需要在概念上锚定这些意象（评委观察到词对正文偏好“命名驱动”而非“问题驱动”的结构），进一步偏离了开放式任务的评价准则。

\textbf{机制三：清洗的信息损失。}与思维链不同，词对环节被设计为可剥离的“脚手架”。但正文的组织方式在生成时已与词对深度耦合（“如第 N 组词所示”“从上述意象中获得灵感”是正文的骨架性连接词），清洗协议在删除这些指代时必须改写句子，即使零残留达标，句法上的缝合痕迹与论证链的中断仍然存在——这从清洗后正文更短（压缩比 0.36--0.43）与逻辑性维度的显著受损（$+1.5$ 分分差）中可见一斑。v1 实验中 17/60 的显性残留则说明，当清洗不彻底时评委还会额外惩罚“引用不可见信息”，这是该策略在工程实现上的又一脆弱点。

\subsection{词对提示何时可能有价值}

本研究的结果并非全盘否定发散式提示，而是刻画了其适用边界。\textbf{第一，原创性维度存在真实且显著的正效应}（$+0.31$ 分，$p<0.001$），说明“自由联想素材”确实能把模型推离常规解空间。\textbf{第二，收益集中于概念构建型任务}：T5（人格模型）上 \wordpairs{} 获得 7/20 的胜场（V4.1-Flash 生成），贡献了全部 15 场胜局中的 9 场，是唯一具有实质竞争力的任务格；该类任务的产出核心是“新颖而自洽的分类学隐喻”，恰是随机意象能直接供给的原料。\textbf{第三，强模型从中获益更少}：V4-Pro 生成时各任务胜场均趋近于零。据此可以给出实践建议：若任务以新颖概念为核心价值且对工程细节不敏感（起名、世界观设定、创意写作中的意象拓展），可考虑小规模使用词对式预热，但应将词对数量降至远小于 200（减少预算挤占）、显式约束词对与任务的相关性（降低自干扰）、或干脆把“发散”与“收敛”拆分为两轮独立调用（让收敛阶段从干净上下文开始，规避清洗损失）。若任务要求可落地方案，则应避免该策略。

\subsection{对提示工程与 LLM 评审方法学的启示}

本研究的证据强化了一个更一般的判断：\textbf{“作答前的额外生成”并非免费}。思维链的成功依赖于中间内容与答案之间的逻辑依赖；一旦前置内容与任务只存在启发式关联（词对、情绪刺激\cite{li2023emotion}、自由联想），其收益就必须与目标稀释、自干扰与后处理损失相权衡，且该权衡很可能为负。对方法学而言，本研究亦示范了一个低成本但严格的评估闭环：交叉双评委消除单一评委偏好、确定性哈希平衡消除位置偏差、正则校验的二次清洗消除交付伪影、全量溯源数据支持独立复核——该协议可直接复用于其他提示策略的评估。位置偏差的测量结果（A 位平均 $+0.27$ 分，方向不完全一致）也再次印证了 Zheng 等人关于 LLM 评委偏差可控但必须显式测量的观点\cite{zheng2023judging,shi2025position}。

\section{有效性威胁}\label{sec:threats}

\textbf{LLM 评委的效度。}本研究完全依赖 LLM 盲评，未包含人工评估。尽管交叉双评委设计、与人类偏好一致性的既有证据\cite{zheng2023judging}以及 v1/v2 两版协议的复现性缓解了单评委风险，评委与人类专家在“原创性”这类主观维度上的判断仍可能存在系统差异。特别是，评委可能对工程完整性存在与任务无关的固有偏好，从而压低意象驱动文本的整体分。\textbf{冗长偏差的反向作用。}LLM 评委通常偏好更长回答\cite{zheng2023judging}，本研究中 \wordpairs{} 清洗后正文反而更短，故冗长偏差只会\emph{低估}词对条件的劣势——即真实差距可能更大，方向性结论不受影响。\textbf{模型家族单一。}两个生成模型均来自 DeepSeek 家族（V4.1-Flash 与 V4-Pro）\cite{deepseekv3}，同源模型的指令遵循与生成风格可能共享相同的失效模式；结论向其他家族模型的外推需要进一步验证（v1 使用 Kimi-K3 得到方向一致的结果，提供了初步的跨家族证据）。\textbf{语言与任务单一。}全部任务与输出均为中文，任务均为一次性开放式生成；其他语言、多轮对话或更长篇幅场景下的表现未考察。\textbf{清洗伪影。}尽管清洗残留为 0，无法完全排除“清洗改写本身而非词对策略”造成部分质量损失的可能；将该损失从策略总效应中剥离需要“对 \baseline{} 输出施加同强度改写”的对照组，这是未来工作的方向之一。\textbf{单一前缀形式。}实验仅测试了一种词对前缀（200 组、两个词、无领域约束）；词对数量、关联度约束与呈现形式的变化可能改变结论，本文结论严格限定于该具体形式。

\section{结论}\label{sec:conclusion}

本文对“作答前先输出 200 组灵感词对”这一发散式提示策略进行了双模型、双评委的受控盲评实验。在 240 条生成与 240 对盲评中，\baseline{} 以 93.3\% 的胜率、$+1.72$ 分的整体质量分差（$p<0.001$）全面占优；词对提示仅在原创性维度获得 $+0.31$ 分的显著增益，集中于概念构建型任务，代价是 $1.60\sim2.02$ 倍的生成成本与清洗后正文的信息密度下降。机理分析将其失效归因于指令目标稀释、联想语境自干扰与清洗信息损失三条通路。

本研究的核心信息是：\textbf{“先发散、后收敛”的直觉不能不加检验地移植到提示设计中}。前置发散素材只有在“新颖性本身即任务核心价值、且发散内容与任务相关性受控”的条件下才可能净收益为正。未来工作包括：(1) 引入人工评估与多模型家族交叉验证；(2) 系统变化词对数量、领域约束与生成-收敛两阶段分离方式，绘制该策略的完整设计空间；(3) 将本协议扩展到其他“过程性前置生成”策略（如情绪刺激、角色扮演预热）的统一评估，为提示工程提供更一般的设计准则。
